\section{Quiver Hecke algebras}
\label{Sec : klr}

\subsection{Quantum group and dual canonical basis}
\label{ssec : quantum group and dual canonical basis}

Let $\mathfrak{g}$ be a symmetric Kac--Moody Lie algebra
with Dynkin index set $I$.
Let $P$ and $P^\vee$ be its weight and coweight lattices,
and let $\{\alpha_i\}_{i \in I}$ and $\{h_i\}_{i \in I}$
be the sets of simple roots and simple coroots.
We set
\[
Q \coloneq \bigoplus_{i \in I}\Z\alpha_i, \qquad
Q_+ \coloneq \bigoplus_{i \in I}\Z_{\geq 0}\alpha_i, \qquad
Q_- \coloneq -Q_+.
\]
Let $W$ be the Weyl group of $\mathfrak{g}$,
and let $s_i$ denote the simple reflection corresponding to $i \in I$.

Let $U_q(\mathfrak{g})$ be the quantum group associated with
$\mathfrak{g}$, that is, the $\Q(q)$-algebra generated by
$e_i, f_i$ $(i \in I)$ and $q^h$ $(h \in P^\vee)$
with the usual defining relations.
We denote by $U_q^+(\mathfrak{g})$ and $U_q^-(\mathfrak{g})$
the subalgebras generated by $\{e_i \mid i \in I\}$
and $\{f_i \mid i \in I\}$, respectively.
For $n \in \Z_{\geq 0}$, set
\[
e_i^{(n)} \coloneq \frac{e_i^n}{[n]_{q_i}!}, \qquad
f_i^{(n)} \coloneq \frac{f_i^n}{[n]_{q_i}!}.
\]
For $\beta \in Q_-$, let $U_q^-(\mathfrak{g})_\beta$
be the homogeneous component of $U_q^-(\mathfrak{g})$
of weight $\beta$.

For each $i \in I$, consider Lusztig's braid group automorphism
$T_i : U_q(\mathfrak{g}) \to U_q(\mathfrak{g})$.
(See \cite{Lusbook}.
Note that, we write $v$ in \cite{Lusbook} as $q$
and adopt $T'_{i, -1}$ for $T_i$.)

Let $w \in W$ and fix a reduced expression
$\widetilde{w} = s_{i_1}\cdots s_{i_r}$.
For $1 \leq k \leq r$, set
\[
\beta_k \coloneq s_{i_1}\cdots s_{i_{k - 1}}(\alpha_{i_k})
\]
and, for $c \in \Z_{\geq 0}$, define
\[
F(c\beta_k) \coloneq T_{i_1}\cdots T_{i_{k - 1}}
\bigl(f_{i_k}^{(c)}\bigr).
\]
For $\mathbf{c} = (c_1, \ldots, c_r) \in \Z_{\geq 0}^r$, set
\[
F(\mathbf{c}, \widetilde{w}) \coloneq
F(c_r\beta_r)F(c_{r - 1}\beta_{r - 1})\cdots F(c_1\beta_1).
\]
It is well-known that the $\Q(q)$-linear subspace spanned by
$\{F(\mathbf{c}, \widetilde{w}) \mid
\mathbf{c} \in \Z_{\geq 0}^r\}$
does not depend on the choice of the reduced word $\widetilde{w}$.
We denote the subspace by $U_q^-(w)$.
Then $\{F(\mathbf{c}, \widetilde{w}) \mid
\mathbf{c} \in \Z_{\geq 0}^r\}$
is a $\Q(q)$-basis of $U_q^-(w)$ and we call it a PBW-basis.

There exists Kashiwara's bilinear form $(\ ,\ )_K$
on $U_q^-(\mathfrak{g})$, which is symmetric and non-degenerate
\cite[Section 3.4]{Kas91}.
By \cite[Lemma 1.4.4, Proposition 38.2.3]{Lusbook},
for $\mathbf{c}, \mathbf{c}' \in \Z_{\geq 0}^r$, we have
\[
\bigl(F(\mathbf{c}, \widetilde{w}),
F(\mathbf{c}', \widetilde{w})\bigr)_K = 0
\quad \text{if } \mathbf{c} \neq \mathbf{c}',
\]
and $\bigl(F(\mathbf{c}, \widetilde{w}),
F(\mathbf{c}, \widetilde{w})\bigr)_K$ is non-zero.
Hence, if we set
\[
F^{\mathrm{up}}(\mathbf{c}, \widetilde{w}) \coloneq
\frac{F(\mathbf{c}, \widetilde{w})}
{\bigl(F(\mathbf{c}, \widetilde{w}),
F(\mathbf{c}, \widetilde{w})\bigr)_K},
\]
then $\{F^{\mathrm{up}}(\mathbf{c}, \widetilde{w}) \mid
\mathbf{c} \in \Z_{\geq 0}^r\}$ is the basis of $U_q^-(w)$
dual to the PBW basis with respect to $(\ ,\ )_K$.
We call it the \emph{dual PBW basis}.

For $\mathbf{c}, \mathbf{c}' \in \Z_{\geq 0}^r$,
we write $\mathbf{c}' < \mathbf{c}$ if there exists $p$
with $1 \leq p \leq r$ such that
\[
c'_1 = c_1, \ldots, c'_{p - 1} = c_{p - 1}, \qquad c'_p < c_p.
\]
Let $\overline{\phantom{x}} : U_q^-(\mathfrak{g})
\to U_q^-(\mathfrak{g})$ be a $\Q$-algebra homomorphism
such that $f_i \mapsto f_i$, $q \mapsto q^{-1}$.
Define a $\Q$-linear map $\sigma : U_q^-(\mathfrak{g})
\to U_q^-(\mathfrak{g})$ by
\[
(\sigma(x), y)_K = \overline{(x, \overline{y})_K}
\qquad \text{for } x, y \in U_q^-(\mathfrak{g}).
\]

\begin{Prop}[{\cite[Theorem 4.29]{Kim12}}]
\label{Prop : dual canonical basis}
There exists a unique $\Q(q)$-basis
\[
\{B^{\mathrm{up}}(\mathbf{c}, \widetilde{w}) \mid
\mathbf{c} \in \Z_{\geq 0}^r\}
\]
of $U_q^-(w)$ satisfying
\[
\sigma\bigl(B^{\mathrm{up}}(\mathbf{c}, \widetilde{w})\bigr)
= B^{\mathrm{up}}(\mathbf{c}, \widetilde{w})
\]
and
\[
F^{\mathrm{up}}(\mathbf{c}, \widetilde{w})
= B^{\mathrm{up}}(\mathbf{c}, \widetilde{w})
+ \sum_{\mathbf{c}' < \mathbf{c}}\phi_{\mathbf{c}, \mathbf{c}'}
B^{\mathrm{up}}(\mathbf{c}', \widetilde{w}), \qquad
\phi_{\mathbf{c}, \mathbf{c}'} \in q\Z[q].
\]
This basis is called the \emph{dual canonical basis} of $U_q^-(w)$.
\end{Prop}

Let $U_q^-(\mathfrak{g})_{\Z[q^{\pm 1}]}$
be a $\Z[q^{\pm 1}]$-algebra generated by
$f_i^{(n)}$ $(i \in I, n \in \Z_{\geq 0})$, and
\[
U_q^-(\mathfrak{g})_{\Z[q^{\pm 1}]}^\vee \coloneq
\{x \in U_q^-(\mathfrak{g}) \mid
(x, U_q^-(\mathfrak{g})_{\Z[q^{\pm 1}]})_K
\subset \Z[q^{\pm 1}]\}.
\]
It is a $\Z[q^{\pm 1}]$-subalgebra of $U_q^-(\mathfrak{g})$
and it has a $\Z[q^{\pm 1}]$-basis called upper global basis
\cite{Kas91}.
Let
\[
U_q^-(w)_{\Z[q^{\pm 1}]}^\vee \coloneq
\bigoplus_{\mathbf{c} \in \Z_{\geq 0}^r}
\Z[q^{\pm 1}]B^{\mathrm{up}}(\mathbf{c}, \widetilde{w}).
\]
Then $U_q^-(w)_{\Z[q^{\pm 1}]}^\vee$
is independent of the choice of reduced word
and it is a subalgebra of $U_q^-(\mathfrak{g})_{\Z[q^{\pm 1}]}^\vee$.
Moreover, its basis
$\{B^{\mathrm{up}}(\mathbf{c}, \widetilde{w}) \mid
\mathbf{c} \in \Z_{\geq 0}^r\}$
is a subset of the upper global basis.
For detail, see \cite[$\S$4.7]{Kim12}.

For $1 \leq s \leq r$, we set
\[
\begin{aligned}
s^- &\coloneq
\max(\{t \in [1, s - 1] \mid i_t = i_s\} \cup \{0\}), \\
s^+ &\coloneq
\min(\{t \in [s + 1, r] \mid i_t = i_s\} \cup \{r + 1\}).
\end{aligned}
\]
We also set
\[
s^{\min} \coloneq \min\{t \in [1, r] \mid i_t = i_s\}, \qquad
s^{\max} \coloneq \max\{t \in [1, r] \mid i_t = i_s\}.
\]
For $1 \leq s \leq t \leq r$ satisfying $i_s = i_t$, we define
\[
D_{\widetilde{w}}[s, t] \coloneq
B^{\mathrm{up}}(\varepsilon_s + \varepsilon_{s^+} + \cdots
+ \varepsilon_{t^-} + \varepsilon_t, \widetilde{w}).
\]
When there is no risk of confusion, we simply write $D[s, t]$.

\subsection{Quiver Hecke algebras}
\label{ssec : quiver Hecke algebras}

Let $\mathfrak{g}$ be a symmetric Kac--Moody Lie algebra
with Cartan matrix $A = (a_{ij})_{i, j \in I}$,
and let $\kk$ be a field of characteristic zero.
For $i, j \in I$, choose a polynomial $Q_{ij}(u, v) \in \kk[u, v]$
of the form
\[
Q_{ij}(u, v) =
\begin{cases}
0 & \text{if } i = j, \\
\displaystyle\sum_{\substack{p, q \in \Z_{\geq 0} \\
p + q = -a_{ij}}} t_{i, j; p, q}u^p v^q & \text{if } i \neq j,
\end{cases}
\]
where $t_{i, j; p, q} \in \kk$ satisfy
\[
Q_{ij}(u, v) = Q_{ji}(v, u), \qquad
t_{i, j; -a_{ij}, 0} \in \kk^\times.
\]

Let $\mathfrak{S}_n = \langle s_1, \ldots, s_{n - 1}\rangle$
be the symmetric group, where $s_k$ is the transposition
interchanging $k$ and $k + 1$.
It acts naturally on $I^n$.
For $\beta \in Q_+$ with $|\beta| = n$, set
\[
I^\beta \coloneq
\{\nu = (\nu_1, \ldots, \nu_n) \in I^n \mid
\alpha_{\nu_1} + \cdots + \alpha_{\nu_n} = \beta\}.
\]

\begin{Def}
\label{Def : quiver Hecke algebra}
For $\beta \in Q_+$ with $|\beta| = n$,
the \emph{quiver Hecke algebra} $R(\beta)$ is the $\kk$-algebra
generated by
\[
\{e(\nu) \mid \nu \in I^\beta\}, \qquad
\{x_k \mid 1 \leq k \leq n\}, \qquad
\{\tau_m \mid 1 \leq m \leq n - 1\},
\]
subject to the following relations:
\[
\begin{gathered}
e(\nu)e(\nu') = \delta_{\nu, \nu'}e(\nu), \qquad
\sum_{\nu \in I^\beta}e(\nu) = 1, \\
x_k x_m = x_m x_k, \qquad x_k e(\nu) = e(\nu)x_k, \\
\tau_m e(\nu) = e(s_m(\nu))\tau_m, \qquad
\tau_k\tau_m = \tau_m\tau_k \quad \text{if } |k - m| > 1, \\
\tau_k^2 e(\nu) = Q_{\nu_k, \nu_{k + 1}}(x_k, x_{k + 1})e(\nu),
\end{gathered}
\]
\[
(\tau_k x_m - x_{s_k(m)}\tau_k)e(\nu) =
\begin{cases}
-e(\nu) & \text{if } m = k \text{ and } \nu_k = \nu_{k + 1}, \\
e(\nu) & \text{if } m = k + 1 \text{ and } \nu_k = \nu_{k + 1}, \\
0 & \text{otherwise},
\end{cases}
\]
and
\[
\begin{aligned}
&(\tau_{k + 1}\tau_k\tau_{k + 1}
- \tau_k\tau_{k + 1}\tau_k)e(\nu) \\
&\quad =
\begin{cases}
\displaystyle
\frac{Q_{\nu_k, \nu_{k + 1}}(x_k, x_{k + 1})
- Q_{\nu_k, \nu_{k + 1}}(x_{k + 2}, x_{k + 1})}
{x_k - x_{k + 2}}e(\nu)
& \text{if } \nu_k = \nu_{k + 2}, \\
0 & \text{otherwise}.
\end{cases}
\end{aligned}
\]
The algebra $R(\beta)$ is $\Z$-graded by
\[
\deg e(\nu) = 0, \qquad
\deg(x_k e(\nu)) = 2, \qquad
\deg(\tau_k e(\nu)) = -a_{\nu_k, \nu_{k + 1}}.
\]
\end{Def}

Let $R(\beta)\text{-}\mathrm{mod}$ be the category of finite-dimensional
$R(\beta)$-modules on which the elements $x_k$ act nilpotently,
and let $R(\beta)\text{-}\mathrm{gmod}$ be the category
of finite-dimensional graded $R(\beta)$-modules.
We set
\[
R\text{-}\mathrm{gmod} \coloneq
\bigoplus_{\beta \in Q_+}R(\beta)\text{-}\mathrm{gmod}.
\]
If $M$ is an $R(\beta)$-module, we set
$\operatorname{wt}(M) \coloneq -\beta \in Q_-$
and call it the \emph{weight} of $M$.
For $\beta, \gamma \in Q_+$, define
\[
e(\beta, \gamma) \coloneq
\sum_{\substack{\nu \in I^\beta \\ \nu' \in I^\gamma}}
e(\nu * \nu') \in R(\beta + \gamma),
\]
where $\nu * \nu'$ denotes the concatenation of $\nu$ and $\nu'$.
For $M \in R(\beta)\text{-}\mathrm{gmod}$
and $N \in R(\gamma)\text{-}\mathrm{gmod}$,
their \emph{convolution product} is
\[
M \circ N \coloneq R(\beta + \gamma)e(\beta, \gamma)
\otimes_{R(\beta) \otimes R(\gamma)}(M \otimes N).
\]

Let $\psi$ be the anti-involution of $R(\beta)$
fixing $e(\nu)$, $x_k$, and $\tau_k$.
For $M \in R(\beta)\text{-}\mathrm{gmod}$, its \emph{graded dual} is
\[
M^* \coloneq \operatorname{Hom}_{\kk}(M, \kk), \qquad
(M^*)_d = \operatorname{Hom}_{\kk}(M_{-d}, \kk),
\]
with the $R(\beta)$-action given by $(xf)(m) = f(\psi(x)m)$.
For $s \in \Z$, the grading shift $q^s M$ is defined by
\[
(q^s M)_d = M_{d - s}.
\]
A simple module $M \in R\text{-}\mathrm{gmod}$
is called \emph{self-dual} if $M \simeq M^*$.

Let $K(R\text{-}\mathrm{gmod})$ be the Grothendieck group
of $R\text{-}\mathrm{gmod}$.
The grading shift and convolution product make it
a $\Z[q, q^{-1}]$-algebra via
\[
[qM] = q[M], \qquad [M][N] = [M \circ N].
\]
The quiver Hecke algebra $R$ is called \emph{symmetric}
if $Q_{ij}(u, v)$ is a polynomial in $u - v$ for every $i, j \in I$.
In what follows, we assume that $R$ is symmetric.

\begin{Thm}[{\cite{KL09, Ro08, VV11}}]
\label{Thm : KLR categorification}
There exists a $\Z[q, q^{-1}]$-algebra isomorphism
\[
\operatorname{ch} : K(R\text{-}\mathrm{gmod})
\xrightarrow{\sim} U_q^-(\mathfrak{g})_{\Z[q^{\pm 1}]}^\vee
\]
under which the isomorphism classes of self-dual simple modules
correspond bijectively to the upper global basis
of $U_q^-(\mathfrak{g})_{\Z[q^{\pm 1}]}^\vee$.
\end{Thm}

Let $w \in W$ and fix a reduced word
$\widetilde{w} = s_{i_1}\cdots s_{i_r}$.
For $1 \leq s \leq t \leq r$ satisfying $i_s = i_t$,
since $D[s, t]$ is a member of the upper global basis,
it follows from Theorem~\ref{Thm : KLR categorification}
that there exists a unique self-dual simple module
$M[s, t] \in R\text{-}\mathrm{gmod}$
satisfying $\operatorname{ch}(M[s, t]) = D[s, t]$.

\begin{Def}[{\cite[$\S$2.2]{strata}}]
\label{Def : Cw}
Let $\mathscr{C}_w$ be the smallest monoidal abelian full subcategory
of $R\text{-}\mathrm{gmod}$ satisfying the following properties:
\begin{enumerate}
    \item $\mathscr{C}_w$ is stable under the subquotients, extensions,
    and grading shifts,
    \item $\mathscr{C}_w$ contains $M[s, s]$ for all $1 \leq s \leq r$.
\end{enumerate}
\end{Def}

Then by \cite{GLSqua}, $M \in R\text{-}\mathrm{gmod}$
belongs to $\mathscr{C}_w$ if and only if $\operatorname{ch}(M)$
belongs to $U_q^-(w)_{\Z[q^{\pm 1}]}^\vee$.
Hence, the following holds:

\begin{Prop}
\label{Prop : Cw categorification}
The isomorphism $\operatorname{ch}$ induces
a $\Z[q, q^{-1}]$-algebra isomorphism
\[
K(\mathscr{C}_w) \xrightarrow{\sim} U_q^-(w)_{\Z[q^{\pm 1}]}^\vee
\]
under which the isomorphism classes of self-dual simple modules
correspond bijectively to the dual canonical basis
of $U_q^-(w)_{\Z[q^{\pm 1}]}^\vee$.
\end{Prop}

For each $1 \leq k \leq r$, let $S_k \coloneq M[k, k]$
and we call it a cuspidal module.

\begin{Prop}[{\cite[$\S$2]{TW16}, \cite[Proposition 2.15]{strata}}]
\label{Prop : KLR PBW parameterization}
\begin{enumerate}
\item [(i)]
For $(c_k)_{k \in [1, r]} \in \Z_{\geq 0}^r$,
the head of $S_r^{\circ c_r} \circ S_{r - 1}^{\circ c_{r - 1}}
\circ \cdots \circ S_1^{\circ c_1}$
is a simple module corresponding to
$B^{\mathrm{up}}(\mathbf{c}, \widetilde{w})$ up to grading shift.
Moreover, if we ignore grading, this correspondence forms
a bijection between $\Z_{\geq 0}^r$ and $\Irr\mathscr{C}_w$.
\item[(ii)] The ring
    $K(\mathscr{C}_w)|_{q=1}$ is the polynomial ring
    generated by $\{[S_s] \mid s \in [1, r]\}$. 
    \end{enumerate}
\end{Prop}

For $L \in \Irr\mathscr{C}_w$, the corresponding sequence $(c_k)_k$
is called \emph{PBW parameter} and we write it as $\mathrm{PBW}(L)$.
